Since @@M0@@ is of finite grading, the graded rings associated with the @@M2@@-ring filtration and the degree filtration on @@M1@@ are identical (6.6.7). We then define @@M3@@ by the identity on @@M4@@ and by:
@@M5@@
and extend by additivity. By relation (6.5.3), this is a homomorphism of graded rings, and it is easy to see that it is compatible with the operations of @@M6@@. The inverse morphism @@M7@@ can be defined as follows: @@M8@@, @@M9@@, the @@M10@@-th coefficient @@M11@@ of @@M12@@ does not depend on the way in which @@M13@@ has been lifted to @@M14@@, and we put:
@@M15@@
The homomorphism @@M16@@ is therefore an isomorphism.

The isomorphism @@M17@@ is defined by the completed Chern class @@M18@@. To verify that it does define an isomorphism, it suffices to verify that the homomorphism obtained by passing to the associated graded rings is one, since the filtrations are discrete. We therefore obtain:
@@M19@@
It suffices to show that the composite is the identity. Since the homomorphisms are homomorphisms of graded @@M20@@-algebras, it suffices, by the definition of the @@M21@@-ring filtration on @@M22@@, to show that it is the identity on the @@M23@@, with @@M24@@. We then obtain:
@@M25@@
Now @@M26@@, by (6.6.4); hence:
@@M27@@
